\clearpage
\phantomsection%
\addcontentsline{toc}{chapter}{ABSTRAK}
\thispagestyle{fancy}

\begin{center}
	\large \bfseries \MakeUppercase{Abstrak}\\
	\bigskip
	
	\normalsize \normalfont \justifying \singlespacing
	Perkembangan media sosial telah mendorong munculnya berbagai bentuk konten digital, termasuk meme yang menggabungkan elemen visual dan teks. Di antara berbagai jenis meme, terdapat meme yang mengandung indikasi \textit{self-harm} dan berpotensi memberikan dampak negatif terhadap pengguna. Tantangan utama dalam mengidentifikasi meme \textit{self-harm} terletak pada sifatnya yang implisit, karena makna sering kali muncul dari interaksi antara gambar dan teks, bukan dari satu modalitas saja. Penelitian ini bertujuan mengembangkan model klasifikasi biner berbasis multimodal untuk mendeteksi meme \textit{self-harm} dengan mengintegrasikan  \textit{Contrastive Language-Image Pre-training} (CLIP) sebagai \textit{image encoder} dan \textit{Efficiently Learning an Encoder that Classifies Token Replacements Accurately} (ELECTRA)sebagai \textit{text encoder} melalui skema \textit{intermediate fusion}. Data penelitian berasal dari gabungan data primer dan sekunder, yang kemudian dianotasi menggunakan pendekatan \textit{ensemble pseudo-labeling} berbasis model multimodal. Proses penelitian meliputi \textit{pre-processing}, pengembangan model unimodal dan multimodal, serta evaluasi menggunakan metrik berbasis \textit{confusion matrix}. Hasil penelitian menunjukkan bahwa model multimodal memperoleh \textit{F1-score} sebesar 96,1\%, lebih tinggi dibandingkan model unimodal gambar sebesar 91,5\% dan model unimodal teks sebesar 92,0\%. Namun, pada data yang ambigu, terutama yang mengandung sarkasme, humor gelap, atau konflik antara teks dan visual, model masih mengalami kesulitan dalam menentukan label secara tepat. Secara keseluruhan, hasil ini menunjukkan bahwa pendekatan multimodal lebih efektif dalam memahami konteks dan mendeteksi meme \textit{self-harm} dibandingkan pendekatan unimodal.\par

    \vspace{20pt}
	\raggedright \textbf{Kata Kunci:} \textit{self-harm}, meme, multimodal, CLIP, \textit{ELECTRA}
	
	\vfill
	
\end{center}
\clearpage
